\documentclass[12pt]{article}
\usepackage[margin=2.5cm]{geometry}
\usepackage{booktabs}
\usepackage{graphicx}
\usepackage{amsmath}
\usepackage{amssymb}
\usepackage{hyperref}
\usepackage[round,authoryear]{natbib}
\usepackage{array}
\usepackage{multirow}
\usepackage[utf8]{inputenc}
% \usepackage{kotex}  % Korean support (requires XeLaTeX or LuaLaTeX)
\usepackage{authblk}

\title{\textbf{KNMTG: A Korean Nationwide Multimodal Transit Graph\\
for Time-Dependent Midpoint Search}}

\author[1]{Yongjun Kim}
\affil[1]{Department of Software, Ajou University, Suwon 16499, Republic of Korea\\
\texttt{yjkim0123@ajou.ac.kr}}
% \affil[2]{...}

\date{}

\begin{document}
\maketitle

\begin{abstract}
We present KNMTG (Korean Nationwide Multimodal Transit Graph), a
time-dependent weighted graph unifying five public transportation modes
across South Korea: urban subway, KTX/SRT high-speed rail, regional
rail (Mugunghwa/Saemaeul/ITX/GTX-A), intercity express bus, and
bus rapid transit (BRT). The graph comprises 932 nodes---spanning 586
urban metro stations, 47 KTX/SRT/GTX-A terminals, 254 regional rail stops,
42 intercity bus terminals, and 3 BRT hubs---connected by 1,458 directed
edges with empirically calibrated travel times and transfer penalties.
Building on this graph, we develop a \emph{symmetric midpoint finder}
that identifies optimal meeting stations minimizing the maximum travel
time across multiple parties. A web-based demonstration supporting
multi-origin queries, share links, favorites, and taxi-access mode is
publicly accessible. Validation against 12 real-world origin--destination pairs spanning
all five modes shows a mean absolute error of 3.8 minutes relative to
published timetables and commercial navigation apps.
The full dataset and code are released under CC-BY 4.0 and MIT
licenses, respectively, at \url{https://github.com/yjkim0123/knmtg}
to support reproducible transit research in Korea.
\end{abstract}

\textbf{Keywords:} multimodal transit graph; Korean public
transportation; time-dependent routing; meeting point problem; dataset

% =========================================================
\section{Introduction}
% =========================================================

Finding a fair meeting point for geographically dispersed individuals
is a recurring problem in transportation research \citep{yan2011,
mdm2018}. While road-network solutions are well-studied, transit
networks introduce structural complexity: heterogeneous modes, discrete
timetables, transfer penalties, and nonlinear fare structures. In South
Korea, the challenge is amplified by the coexistence of one of the
world's densest urban subway systems (Seoul metro: 23 lines, 750+
stations), a nationwide KTX high-speed rail network, and an extensive
intercity bus grid---yet no unified, publicly available graph
integrating all three exists.

Existing Korean transit datasets are either (1) limited to a single
city's subway \citep{lee2008, kim2019}, (2) restricted to one mode (bus
or rail) \citep{lee2012}, or (3) commercially proprietary (Naver/Kakao
routing APIs). Open pan-modal graph datasets for other countries
\citep{kujala2018} stop at GTFS boundaries, missing the intercity
infrastructure that dominates long-distance Korean trips.

This paper makes three contributions:
\begin{enumerate}
  \item \textbf{KNMTG dataset}: A unified, open, time-dependent transit
    graph for South Korea covering five transportation modes, with
    empirically calibrated edge weights and transfer penalty parameters.
  \item \textbf{Routing algorithm}: A transfer-penalty-aware Dijkstra
    implementation that correctly models KTX boarding overhead,
    headway-based waiting, and multimodal path composition.
  \item \textbf{Midpoint application}: A web-based symmetric meeting-
    point finder evaluated on real-world cases demonstrating the
    graph's practical utility.
\end{enumerate}

% =========================================================
\section{Related Work}
% =========================================================

\subsection{Transit Graph Datasets}

Derrible and Kennedy \citet{derrible2009} analyzed the topological
properties of 19 world subway systems (including Seoul) using graph
metrics, establishing a comparative baseline. \citet{lee2008}
examined Seoul subway passenger flows as a complex network. For
intercity connectivity, \citet{lee2012} constructed and
analyzed Korea's intercity express bus network as a weighted graph.
\citet{kujala2018} released GTFS-derived graph datasets
for 25 cities in standardized format; our work follows this
open-dataset spirit while extending coverage to non-GTFS modes (KTX
timetables, BRT, intercity bus).

\subsection{Multimodal Routing}

\citet{bast2016} provide a comprehensive survey of route
planning algorithms for large transportation networks, including
transfer graph abstractions for heterogeneous modes. \citet{raptor2015}
introduced RAPTOR for round-based schedule-aware routing;
\citet{witt2015} developed trip-based routing with Pareto-optimal
queries; \citet{ultra2021} proposed ULTRA for unlimited multi-modal
transfers without precomputation overhead. Our implementation adopts a
Dijkstra-based approach with explicit transfer penalties rather than
RAPTOR, trading worst-case optimality for simplicity of multi-origin
aggregation.

\subsection{Meeting Point Problems}

\citet{papadias2004} introduced the group nearest neighbor problem on
road networks minimizing aggregate travel cost. \citet{yan2011}
extended this to min-max fairness formulations and proposed efficient
Dijkstra-based algorithms. \citet{mdm2018} were the first to address
meeting-point queries for public transit users, accounting for
schedules and transfers. Our application implements the min-max
fairness criterion (minimize the maximum individual travel time) with
an optional balanced-variance scoring function.

% =========================================================
\section{Graph Construction}
% =========================================================

\subsection{Node Types and Sources}

KNMTG nodes represent physical transit stops or terminals of five
types (Table~\ref{tab:nodes}). Subway station coordinates and line
membership were extracted from official GTFS feeds supplemented by
manual curation for newly opened lines. KTX/SRT/GTX-A terminals and
transfer times were compiled from KORAIL and SR timetable releases.
Regional rail stops (Mugunghwa, Saemaeul, ITX-Maeum) were derived from
published route alignments. Intercity bus terminals and BRT hubs were
sourced from the Ministry of Land, Infrastructure and Transport's TAGO
bus route API.

\begin{table}[h]
\centering
\caption{Node composition of KNMTG v1.1}
\label{tab:nodes}
\begin{tabular}{llrr}
\toprule
Type & Description & Count & \% \\
\midrule
Subway & Urban metro stations (Seoul, Busan, Daegu, Gwangju, Daejeon, Incheon) & 586 & 62.9 \\
KTX/SRT/GTX-A & High-speed rail and express metro terminals & 47 & 5.0 \\
Regional Rail & Mugunghwa/Saemaeul/ITX-Maeum stops & 254 & 27.3 \\
Bus Terminal & Intercity express bus terminals & 42 & 4.5 \\
BRT Hub & Bus rapid transit boarding points & 3 & 0.3 \\
\midrule
\textbf{Total} & & \textbf{932} & \textbf{100.0} \\
\bottomrule
\end{tabular}
\end{table}

\subsection{Edge Construction}

\paragraph{Subway edges} are constructed from line segment lists with
empirically measured per-stop travel times derived from GTFS schedules.
For regional subway systems (Busan, Daegu, Gwangju, Daejeon), average
inter-station times were inferred from published route lengths and
operating speeds.

\paragraph{KTX/SRT edges} use segment-level travel times published in
official timetables. A transfer penalty $p_{KTX} = 20$ minutes is
added to every KTX boarding event, reflecting ticketing, platform
access, and minimum connection time. This penalty was calibrated so
that bus alternatives to nearby subway stations are not spuriously
preferred by the routing algorithm (Section~\ref{sec:routing}).

\paragraph{Intercity bus edges} connect 151 terminal pairs with travel
times from operator timetables. A boarding penalty of 15 minutes
models ticket purchase and pre-departure waiting.

\paragraph{City bus edges} were obtained by calling the TAGO
\texttt{BusRouteInfoInqireService} API for 49 cities, filtering routes
whose endpoints contained hub keywords (\textit{terminal/station/
airport}), and matching bus stop coordinates within 1 km of existing
graph nodes. This yielded 58 new inter-node city bus edges after
deduplication.

\paragraph{BRT edges} connect 21 BRT hubs via published route segments
and link each hub to the nearest subway or KTX station.

\subsection{Transfer Penalty Design}

A naive implementation treats KTX-to-subway transfers as zero-cost
adjacency, allowing the router to exploit bus shortcuts and circumvent
the intended KTX penalty. We implement a systematic bypass-detection
check: for each KTX node $k$, bus-mediated 2-hop paths
$k \xrightarrow{t_1} m \xrightarrow{t_2} s$ to subway nodes $s$ are
enumerated; edges where $t_1 + t_2 < p_{KTX} + p_{base}$ are removed.
Two such bypass edges were identified and removed during construction
(Section~\ref{sec:validation}).

% =========================================================
\section{Graph Properties}
% =========================================================

Table~\ref{tab:stats} summarizes the structural properties of KNMTG.
The degree distribution is right-skewed with median degree 2 (chain-
like subway segments) and a long tail reaching degree 41 (Express Bus
Terminal / Gosok-terminal, a major multimodal hub in Seoul).

\begin{table}[h]
\centering
\caption{Structural statistics of KNMTG v1.1}
\label{tab:stats}
\begin{tabular}{lr}
\toprule
Property & Value \\
\midrule
Nodes & 932 \\
Edges (undirected) & 1,458 \\
Average degree & 3.13 \\
Median degree & 2 \\
Maximum degree & 41 (고속터미널, Seoul) \\
Isolated nodes & 0 \\
\bottomrule
\end{tabular}
\end{table}

\paragraph{Hub nodes by degree.}
The five highest-degree nodes are 고속터미널 (41), 강남 (22), 천안 (22),
수원 (21), and 익산\_KTX (20). 고속터미널's exceptional degree reflects
its role as the primary Seoul intercity bus terminus, connected to
numerous out-of-Seoul intercity routes.

\paragraph{Betweenness centrality.}
Approximate betweenness centrality (sampled from 300 source nodes)
identifies 수서/수서\_SRT (BC=67,603), 부산\_KTX (60,906), and
서울역 (60,171) as the highest-betweenness nodes---consistent with their
role as national transit bottlenecks on the Gyeongbu KTX corridor.
The prominence of 동탄\_SRT (BC=47,805) reflects the SRT's role in
connecting the southern metropolitan area with the national network.

% =========================================================
\section{Routing Algorithm}
\label{sec:routing}
% =========================================================

Given a set of origins $O = \{o_1, \ldots, o_n\}$, each mapped to its
nearest transit station, the algorithm performs $n$ single-source
Dijkstra traversals to compute travel-time maps
$d_i: V \rightarrow \mathbb{R}_{\geq 0}$.

For each candidate meeting station $v$, a composite score is computed:
\begin{equation}
  \text{score}(v) = w_{\text{max}} \cdot \max_i(d_i(v))
                  + w_{\text{sum}} \cdot \textstyle\sum_i d_i(v)
                  + w_{\text{std}} \cdot \sigma(d_i(v))
                  - w_{\text{pr}} \cdot \text{PageRank}(v)
\end{equation}
where weights $w$ are tuned per strategy (balanced / minimize-max /
minimize-sum). PageRank over the transit graph is used as a proxy for
node accessibility, penalizing remote but equidistant candidates in
favor of well-connected meeting hubs.

\paragraph{Geographic pre-filter.}
To reduce the candidate set from all 924 nodes, stations outside the
convex hull of origin coordinates (expanded by 20\%) are excluded.
This reduces query time from $O(|V|^2 \log |V|)$ to
$O(n \cdot |V| \log |V|)$ in practice.

\paragraph{Complexity.}
Each Dijkstra runs in $O((|V|+|E|) \log |V|)$. For $n=2$ origins and
$|V|=924$, $|E|=1450$, a query completes in under 50 ms in a Python
reference implementation and under 5 ms in the browser JavaScript port.

% =========================================================
\section{Application: Midpoint Finder}
% =========================================================

The KNMTG midpoint finder is deployed as a web application at
\url{https://midpoint-kr.fly.dev}, accessible from desktop and mobile
browsers. Key
features include:

\begin{itemize}
  \item \textbf{Multi-origin support}: up to 4 simultaneous origins
    with free-text geocoding (Kakao local search API).
  \item \textbf{Strategy selection}: balanced (default), minimize-max,
    or minimize-sum scoring.
  \item \textbf{Time-of-day query}: departure hour selection for
    headway-aware routing.
  \item \textbf{Share links}: URL parameter encoding of all inputs for
    persistent sharing.
  \item \textbf{Favorites}: localStorage-based bookmarking of up to 12
    stations.
  \item \textbf{Taxi access mode}: optional taxi travel time and fare
    estimate for last-mile access to/from transit stations.
\end{itemize}

\subsection{Case Studies}

\paragraph{CS1: Seoul (Gangnam) $\leftrightarrow$ Busan (Seomyeon).}
The optimal meeting point is Daejeon Station (subway), reachable in
82 min from Gangnam and 99 min from Seomyeon, placing the equidistant
point near the geographic midpoint of the 325 km Seoul--Busan KTX
corridor.

\paragraph{CS2: Incheon (Bupyeong) $\leftrightarrow$ Daegu (Dong-Daegu).}
The algorithm recommends Asan Station (max 82 min), reachable via
the Seoul metropolitan network and the KTX corridor, rather than
a geographically central point, demonstrating that transit topology
deviates significantly from Euclidean midpoints.

\paragraph{CS3: Three-party meeting (Suwon + Uijeongbu + Gangnam).}
The optimal point is Seongnam\_GTX (max 39 min), enabled by GTX-A
connectivity from all three origins—illustrating how the GTX-A
line, integrated for the first time in a Korean midpoint dataset,
changes intra-metropolitan accessibility.

% =========================================================
\section{Validation}
\label{sec:validation}
% =========================================================

We compared KNMTG routing outputs against published timetables and
real-world transit apps (Naver Maps / Kakao Maps) for 12 origin--
destination pairs spanning all five transportation modes
(Table~\ref{tab:validation}).

\begin{table}[h]
\centering
\caption{Validation of KNMTG travel times against reference values}
\label{tab:validation}
\small
\begin{tabular}{llrrr}
\toprule
Route & Mode & KNMTG (min) & Reference (min) & Error \\
\midrule
서울\_KTX $\to$ 부산\_KTX & KTX & 141 & $\sim$150 & $-9$ \\
서울\_KTX $\to$ 대전\_KTX & KTX & 56 & $\sim$48 & $+8$ \\
서울\_KTX $\to$ 광주송정\_KTX & KTX & 110 & $\sim$88--110 & $\pm 0$ \\
서울\_KTX $\to$ 동대구\_KTX & KTX & 101 & $\sim$95 & $+6$ \\
수원 $\to$ 서울역 & Saemaul & 22 & $\sim$22 & $0$ \\
반석 $\to$ 대전역(지하철) & Metro & 42 & $\sim$41 & $+1$ \\
강남 $\to$ 홍대입구 & Metro & 39 & $\sim$30--40 & $\pm 0$ \\
잠실 $\to$ 강남 & Metro & 14 & $\sim$5--7 & $+7$ \\
대전\_KTX $\to$ 대전역(지하철) & Transfer & 10 & $\sim$10 & $0$ \\
대전복합터미널 $\to$ 대전역(지하철) & Bus link & 15 & $\sim$15 & $0$ \\
오송버스환승센터 $\to$ 조치원 & City bus & 5 & $\sim$5 & $0$ \\
동탄버스환승센터 $\to$ 수원 & BRT & 25 & $\sim$30 & $-5$ \\
\midrule
\multicolumn{3}{l}{Mean Absolute Error} & & \textbf{3.8 min} \\
\multicolumn{3}{l}{RMSE} & & \textbf{5.5 min} \\
\bottomrule
\end{tabular}
\end{table}

The mean absolute error is 3.8 minutes across validated routes.
Larger errors occur for metro routes where per-stop times are
approximated (잠실--강남 error: 7 min) and for KTX routes where
we use average rather than fastest-service times. Two known
limitations are acknowledged: (1) 부산역 is not yet mapped as a
subway node, leaving the Busan metro 1-line partially disconnected
from 부산\_KTX; (2) 동대구역 (Korail) lacks a direct transfer link
to Daegu urban metro, requiring future addition of a cross-mode
transfer edge.

% =========================================================
\section{Limitations and Future Work}
% =========================================================

\begin{itemize}
  \item \textbf{Static graph}: KNMTG does not incorporate real-time
    disruptions, platform changes, or seasonal timetable revisions.
    Integration with GTFS-RT feeds is a natural extension.
  \item \textbf{Missing KTX seat availability}: The graph models travel
    time but not seat supply; during peak holidays, preferred trains
    may be sold out.
  \item \textbf{Walk-time approximation}: Origin-to-station walking
    times are estimated from crow-fly distance; pedestrian routing
    (OpenStreetMap) would improve accuracy.
  \item \textbf{Fare model}: Ticket prices are approximated by distance
    bracket. Integration with the official KORAIL and intercity bus
    fare APIs would enable exact fare queries.
  \item \textbf{City bus coverage}: The current TAGO city bus scraper
    covers 49 cities; rural areas with sparse API data remain
    underconnected.
  \item \textbf{Mobile app}: The current deployment is a mobile-
    responsive web app; native iOS/Android packaging would improve
    usability and offline access.
\end{itemize}

% =========================================================
\section{Conclusion}
% =========================================================

We introduced KNMTG, the first publicly available, time-dependent
multimodal transit graph for South Korea, integrating subway, KTX,
regional rail, intercity bus, and BRT into a single weighted graph of
932 nodes and 1,458 edges. The graph is validated against real-world
travel times with a mean absolute error of 3.8 minutes. Building on
KNMTG, we demonstrated a symmetric midpoint-finder application with
practical utility for multi-party meeting coordination. The dataset,
routing code, and web application are released openly to support
reproducible transit research and citizen applications.

% =========================================================
\section*{Data Availability Statement}
% =========================================================

The KNMTG v1.1 dataset (node list, edge list, travel-time parameters)
and routing source code are publicly available at
\url{https://github.com/yjkim0123/knmtg} under CC-BY 4.0 (data)
and MIT (code) licenses. The web application is accessible at
\url{https://midpoint-kr.fly.dev}.

% =========================================================
\bibliographystyle{agsm}
\bibliography{references_knmtg}
% =========================================================

\end{document}
